% دسته‌بندی روش‌های برآورد و مهندسی گاف
% سبک مشترک نمودارهای سند پایه‌ها (فقط داخل محیط latin فراخوانی شود)
\tikzset{
	qafbox/.style={
		draw=black!55,
		rounded corners=2pt,
		align=center,
		inner sep=3pt,
		thick,
		fill=white,
		font=\footnotesize\linespread{1}\selectfont,
		minimum height=0.95cm
	},
	qafwide/.style={
		qafbox,
		text width=2.55cm
	},
	qafnarr/.style={
		qafbox,
		text width=2.05cm
	},
	qaftitle/.style={
		font=\small\bfseries\linespread{1}\selectfont,
		align=center
	},
	qafarr/.style={
		-{Latex[length=2.2mm]},
		thick,
		draw=black!65
	},
	qafpanel/.style={
		draw,
		thick,
		rounded corners=4pt,
		inner xsep=8pt,
		inner ysep=8pt
	}
}

\begin{tikzpicture}[
	x=1cm,
	y=1cm,
	font=\footnotesize,
	gbox/.style={qafbox, text width=2.35cm, minimum height=1.05cm}
	]
	\node[gbox, fill=blue!12, text width=3.4cm] (why) at (0,3.4) {Why the gap?};
	\node[gbox, fill=green!12] (cl) at (-5.2,1.6) {Classical spectrum};
	\node[gbox, fill=orange!14] (dyn) at (-1.75,1.6) {Dynamics / spectroscopy};
	\node[gbox, fill=cyan!14] (var) at (1.75,1.6) {Variational / subspace};
	\node[gbox, fill=red!12] (ft) at (5.2,1.6) {FTQC / QPE};
	\node[gbox, fill=green!10, text width=2.55cm] (cl2) at (-5.2,0) {Lanczos, DMRG, QMC};
	\node[gbox, fill=orange!10, text width=2.55cm] (dyn2) at (-1.75,0) {Oscillation, quench, $C(\tau)$};
	\node[gbox, fill=cyan!10, text width=2.55cm] (var2) at (1.75,0) {VQD, QSE, Krylov};
	\node[gbox, fill=red!10, text width=2.55cm] (ft2) at (5.2,0) {Phase est., counting};
	\node[gbox, fill=violet!12, text width=4.6cm] (eng) at (0,-1.7) {Gap engineering / shortcuts};
	\node[font=\scriptsize, text=black!60] at (0,-2.45) {schedule, catalyst, CD, encoding};
	\draw[qafarr] (why) -- (cl);
	\draw[qafarr] (why) -- (dyn);
	\draw[qafarr] (why) -- (var);
	\draw[qafarr] (why) -- (ft);
	\draw[qafarr] (cl) -- (cl2);
	\draw[qafarr] (dyn) -- (dyn2);
	\draw[qafarr] (var) -- (var2);
	\draw[qafarr] (ft) -- (ft2);
	\draw[qafarr] (cl2.south) -- ++(0,-0.35) -| (eng.north);
	\draw[qafarr] (dyn2.south) -- ++(0,-0.35) -| (eng.north);
	\draw[qafarr] (var2.south) -- ++(0,-0.35) -| (eng.north);
	\draw[qafarr] (ft2.south) -- ++(0,-0.35) -| (eng.north);
\end{tikzpicture}
